% ============================================================================
% LANGENTWURF LE-003: Das Apple-Ökosystem
% Profil: S (Senioren 65+) | Tier: 2 (Standard)
% ============================================================================

\documentclass[11pt,a4paper]{article}
\usepackage[utf8]{inputenc}
\usepackage[T1]{fontenc}
\usepackage[ngerman]{babel}
\usepackage{geometry}
\usepackage{fancyhdr}
\usepackage{lastpage}
\usepackage{xcolor}
\usepackage{tcolorbox}
\usepackage{enumitem}
\usepackage{longtable}
\usepackage{booktabs}
\usepackage{hyperref}
\usepackage{amssymb}

\geometry{left=2.5cm, right=2.5cm, top=2.5cm, bottom=2.5cm}
\definecolor{fach}{HTML}{b35c00}
\definecolor{gray}{HTML}{666666}
\definecolor{lightgray}{HTML}{f5f5f5}
\definecolor{successgreen}{HTML}{2a7a4b}

\tcbuselibrary{skins,breakable}
\newtcolorbox{infobox}[1][]{ colback=lightgray, colframe=fach, fonttitle=\bfseries, title=#1, breakable }
\newtcolorbox{scorebox}{ colback=white, colframe=fach, fonttitle=\bfseries, title=Didaktik-Score, breakable }

\pagestyle{fancy}
\fancyhf{}
\fancyhead[L]{\small\textcolor{gray}{Digitale Grundlagen -- 003 Apple-Ökosystem}}
\fancyhead[R]{\small\textcolor{gray}{LE Tier 2 / Profil S}}
\fancyfoot[C]{\small\thepage{} / \pageref{LastPage}}
\renewcommand{\headrulewidth}{0.4pt}

\begin{document}

\begin{titlepage}
    \centering
    \vspace*{2cm}
    {\Large\textcolor{gray}{Digitale Grundlagen -- Senioren 65+}}\\[2cm]
    {\Huge\textcolor{fach}{\textbf{Langentwurf}}}\\[0.5cm]
    {\large Tier 2 | Profil S}\\[2cm]
    {\LARGE\textbf{003 -- Das Apple-Ökosystem}}\\[0.5cm]
    {\large Modul A: Geräte \& Zugänge}\\[2cm]
    \begin{tabular}{rl}
        \textbf{Zeitrahmen:} & 60--75 Minuten \\
        \textbf{Vorkenntnisse:} & Einheit 001 (Apple-ID kennen) \\
    \end{tabular}
    \vfill
    {\normalsize Stand: \today}
\end{titlepage}

\tableofcontents
\newpage

\section{Bedingungsanalyse}

\subsection{Ausgangssituation}
\begin{itemize}
    \item Teilnehmer kennen die Apple-ID aus Einheit 001
    \item Verstehen noch nicht, was alles damit verbunden ist
    \item Begriffe wie ``iCloud'', ``App Store'', ``Wo ist?'' sind unklar
    \item Wissen nicht, dass Fotos automatisch synchronisiert werden
\end{itemize}

\subsection{Typische Verwirrung}
\begin{itemize}
    \item ``Was ist iCloud? Brauche ich das?''
    \item ``Warum erscheint mein Foto plötzlich auf dem iPad?''
    \item ``Muss ich Apps auf jedem Gerät neu kaufen?''
    \item ``Kann jemand meine Fotos sehen?''
\end{itemize}

\section{Sachanalyse}

\subsection{Kernkonzepte}
\begin{enumerate}
    \item \textbf{Apple-ID als Zentrum:} Alle Dienste sind mit einer ID verbunden
    \item \textbf{iCloud:} Speicher in der ``Wolke'' (Fotos, Kontakte, Backups)
    \item \textbf{App Store:} Apps kaufen/laden, einmal kaufen = überall nutzen
    \item \textbf{``Wo ist?'':} Geräte orten, sperren, löschen
    \item \textbf{iMessage/FaceTime:} Kostenlose Kommunikation zwischen Apple-Geräten
    \item \textbf{Synchronisation:} Änderungen erscheinen automatisch überall
\end{enumerate}

\subsection{Alltagsanalogien}
\begin{tabular}{|l|p{8cm}|}
\hline
\textbf{Dienst} & \textbf{Analogie} \\
\hline
Apple-ID & Gästekarte in einer Hotelkette \\
iCloud & Schließfach, das von überall erreichbar ist \\
App Store & Laden, in dem du einmal kaufst und überall mitnimmst \\
``Wo ist?'' & GPS-Tracker für deine Geräte \\
Synchronisation & Automatische Kopiermaschine \\
\hline
\end{tabular}

\section{Didaktische Analyse}

\subsection{Stellung in der Reihe}
Einheit 003 vertieft das Verständnis der Apple-ID aus Einheit 001. Hier geht es um das ``große Ganze'' -- wie alles zusammenhängt.

\subsection{Legitimation}
\begin{itemize}
    \item Verständnis der Synchronisation verhindert Verwirrung
    \item ``Wo ist?'' ist sicherheitsrelevant (Geräteverlust)
    \item Wissen über iCloud-Speicher verhindert Fehlkäufe
\end{itemize}

\subsection{Didaktische Reduktion}
\textbf{Ausgeklammert:} Familienfreigabe, iCloud+, Apple One, detaillierte iCloud-Einstellungen.

\textbf{Fokus:} Grundverständnis der Hauptdienste und ihrer Verbindung.

\section{Lernziele}

\begin{tabular}{|c|p{7cm}|c|c|}
\hline
\textbf{Nr.} & \textbf{Die Teilnehmer können...} & \textbf{Stufe} & \textbf{Niveau} \\
\hline
FZ1 & erklären, dass die Apple-ID alle Dienste verbindet & Verstehen & Basis \\
FZ2 & die Hauptdienste (iCloud, App Store, ``Wo ist?'') benennen & Wissen & Basis \\
FZ3 & verstehen, warum Fotos auf mehreren Geräten erscheinen & Verstehen & Basis \\
FZ4 & ``Wo ist?'' auf ihrem Gerät finden & Anwenden & Basis \\
FZ5 & ihren iCloud-Speicherstand prüfen & Anwenden & Vertiefung \\
\hline
\end{tabular}

\section{Methodische Analyse}

\subsection{Vorgehen}
\begin{enumerate}
    \item \textbf{Einstieg:} ``Warum ist das Foto auf beiden Geräten?'' -- Rätsel
    \item \textbf{Überblick:} Diagramm zeigen (Apple-ID in der Mitte, Dienste drum herum)
    \item \textbf{Einzelne Dienste:} Jeweils kurz erklären mit Analogie
    \item \textbf{Demonstration:} ``Wo ist?'' und iCloud-Speicher am Gerät zeigen
    \item \textbf{Übungen:} Zuordnung (Situation $\rightarrow$ Dienst)
\end{enumerate}

\subsection{Material}
Das vorhandene Material enthält:
\begin{itemize}
    \item Seite 1: Grundidee + Diagramm (Apple-ID zentral)
    \item Seite 2: Dienste im Detail (iCloud, App Store, ``Wo ist?'', iMessage)
    \item Seite 3: Übungen + Checkliste
\end{itemize}

\section{Verlaufsplanung}

\begin{longtable}{|p{1cm}|p{2cm}|p{5cm}|p{3cm}|p{2cm}|}
\hline
\textbf{Zeit} & \textbf{Phase} & \textbf{Inhalt} & \textbf{TN-Aktivität} & \textbf{Material} \\
\hline
0--10 & Einstieg & ``Das Foto-Rätsel'': Warum erscheint ein Foto auf mehreren Geräten? & Nachdenken, Vermuten & -- \\
\hline
10--25 & Überblick & Apple-ID als Zentrum, Diagramm erklären, Hotelgast-Analogie & Zuhören, Diagramm studieren & S.1 \\
\hline
25--45 & Dienste & iCloud, App Store, ``Wo ist?'', iMessage einzeln erklären & Mitmachen am Gerät & S.2 \\
\hline
45--55 & Übungen & Zuordnungsaufgaben, Richtig/Falsch & Bearbeiten & S.3 \\
\hline
55--65 & Sicherung & Checkliste, iCloud-Speicher prüfen & Abhaken, Notieren & S.3 \\
\hline
\end{longtable}

\section{Lösungen}

\textbf{Übung: Was nutze ich wofür?}
\begin{enumerate}
    \item Foto auf iPad sehen: \textbf{iCloud}
    \item Neue App laden: \textbf{App Store}
    \item iPhone verloren: \textbf{``Wo ist?''}
    \item Videotelefonieren: \textbf{FaceTime}
    \item Fotos sichern: \textbf{iCloud}
    \item Kostenlose Nachricht: \textbf{iMessage}
\end{enumerate}

\textbf{Richtig/Falsch:}
1. F, 2. R, 3. F, 4. R, 5. F, 6. R

\section{Didaktik-Score}

\begin{scorebox}
\begin{tabular}{|c|l|c|c|p{3.5cm}|}
\hline
\# & Dimension & Gew. & Score & Notiz \\
\hline
1 & Differenzierung & 15\% & 8/10 & Basis/Vertiefung erkennbar \\
2 & Taxonomie & 10\% & 9/10 & Verstehen dominiert (passend) \\
3 & Alltagsrelevanz & 10\% & 10/10 & Foto-Rätsel sehr relevant \\
4 & Aktivierung & 10\% & 8/10 & Übungen + Gerät-Demo \\
5 & Scaffolding & 10\% & 9/10 & Gutes Diagramm, Analogien \\
6 & Feedback & 10\% & 8/10 & Lösungen vorhanden \\
7 & Sprachliche Klarheit & 10\% & 9/10 & Analogien helfen \\
8 & Motivation & 10\% & 10/10 & Rätsel-Einstieg motiviert \\
9 & Barrierefreiheit & 10\% & 9/10 & Gute Struktur, Farben \\
10 & Praktikabilität & 5\% & 8/10 & 65 Min evtl. knapp \\
\hline
\multicolumn{3}{|r|}{\textbf{GESAMT:}} & \textbf{8.95/10} & \textcolor{successgreen}{Knapp unter Freigabe} \\
\hline
\end{tabular}
\end{scorebox}

\subsection{Verbesserungsvorschlag}
\textbf{Praktikabilität:} Bei Zeitmangel ``Wo ist?'' nur zeigen, nicht selbst ausprobieren lassen.

\end{document}
